\section{Measures and models for rational fibrations}
\label{sec:models}

The induction will transfer two metrics along a rational map. Divisor
orders are best controlled on a normal equidimensional model; integration
requires a smooth model. We first record how to retain both models at
once. We also specify the analytic conventions behind the finite-volume
hypothesis.

\subsection{Adjoint measures and extension of weights}

In a holomorphic frame $v$ of a line bundle $J$, a semipositive singular
metric has squared frame norm $|v|^2_h=e^{-\phi}$ with $\phi$
plurisubharmonic (psh) and not identically $-\infty$. In particular
$\phi$ is locally bounded above and $|v|_h^2$ has a positive local
lower bound. A metric with locally bounded weights has bounds on both
sides. Pullbacks and products of the metrics used below preserve
semipositivity.

For $J=-K_Z+D$, write the $J$-valued top form $s_D$ in coordinates as
\[
 s_D=g(z)\,dz_1\wedge\cdots\wedge dz_d\otimes v.
\]
Its squared norm is the measure
\begin{equation}\label{eq:adjoint-measure}
       |s_D|_h^2=|g(z)|^2e^{-\phi(z)}d\lambda_z,
\end{equation}
with one fixed dimensional normalization of Lebesgue measure. The
expression is independent of coordinates and frame. Its total mass
need not be finite merely because the curvature is semipositive.
For a birational morphism $\pi:Y\to Z$ between smooth varieties,
the pulled-back form includes the Jacobian:
\[
 \pi^*s_D\in H^0(Y,K_Y+\pi^*J),\qquad
 \operatorname{div}(\pi^*s_D)=K_{Y/Z}+\pi^*D.
\]
Here $K_{Y/Z}=K_Y-\pi^*K_Z$ is the effective exceptional Jacobian
divisor. Change of variables preserves the total mass in
\eqref{eq:adjoint-measure}. All assertions about singular metrics and
fiber integrals will be made as assertions about measures or almost
everywhere defined functions, followed by their psh representatives.

We use two removable-singularity facts. A psh function on the
complement of an analytic set extends if it is locally bounded above.
Across an analytic set of codimension at least two in a smooth space,
this upper bound is automatic. To see the latter locally, project a
small polydisc in a generic complex direction. Boundaries of suitable
parallel discs lie in a compact set outside the analytic subset; almost
all such discs miss that subset. The maximum principle gives an upper
bound on those discs, and the submean inequality gives the same bound
at every point of the complement. Extension by upper limits is then
psh. These statements apply in line-bundle frames, and a uniform lower
bound survives extension by upper limits.

\subsection{A normal model and a smooth resolution}

\begin{lemma}[Prepared model]\label{lem:prepared-model}
Let $Z$ be smooth connected projective and $S_0$ an integral complex
variety. Let $Z\dashrightarrow S_0$ be a dominant rational map with connected
general fiber. There is a smooth projective birational model $S$ of
$S_0$ and a diagram
\[
\begin{tikzcd}[row sep=large,column sep=large]
Y\arrow[r,"r"]\arrow[dr,"f"']\arrow[rr,bend left=22,"\pi"]
 &\bar Y\arrow[r,"\bar\pi"]\arrow[d,"\bar f"]&Z\\
 &S&
\end{tikzcd}
\]
in which $\bar Y$ is normal projective, $Y$ is smooth projective,
$r$ and $\bar\pi$ are birational morphisms, and $\bar f$ is
surjective and equidimensional. The construction may start above a
prescribed birational model of $Z$. If an algebraic torus acts on $Z$
and the rational map is invariant, the diagram can be chosen
equivariant with trivial action on $S$.

Every prime divisor on $\bar Y$ either dominates $S$ or maps onto a
prime divisor of $S$. At a general point of the strict transform in
$Y$ of a prime $Q\subset\bar Y$ above a base prime $T$, there are
coordinates in which
\begin{equation}\label{eq:power-map}
 f(t,x_2,\ldots,x_d)=(t^e,x_2,\ldots,x_s),
 \qquad e=\ord_Q(\bar f^*T)>0,
\end{equation}
where $s=\dim S$ and $d=\dim Z$.
\end{lemma}

\begin{proof}
First replace $S_0$ birationally by a projective model of its function
field and compose the rational map with that replacement. This preserves
the generic fiber and any trivial base action.
Resolve the graph of the resulting rational map, starting above the prescribed
model if one has been given. Flatten the resulting projective morphism
by a modification of its base, using the strict transform
\cite{RaynaudGruson1971}. Resolve the modified base and take the main
component of the base change, namely the closure of the component over
the common dense open. Before taking that component, the flat family
has all fibers of the generic dimension $d-s$; base change preserves
this. A closed component has fiber dimension at most $d-s$.
Normalize the main component to obtain $\bar Y$. Normalization is
finite, so this upper bound persists. The dimension inequality for a
dominant map gives the reverse bound for every irreducible fiber
component. Thus $\bar f$ is equidimensional. Resolve $\bar Y$ to
obtain $Y$; no equidimensionality is asserted for $f$.

For an invariant rational map, use its invariant graph. The strict
transform and main component are invariant, the action lifts uniquely
through normalization, and characteristic-zero equivariant resolution
supplies the smooth source. All base modifications have trivial torus
action, so the resulting maps are equivariant.

If a prime $Q$ has image of codimension at least two in $S$, its generic
fiber over its image has dimension at least
$(d-1)-(s-2)=d-s+1$, contradicting equidimensionality. Hence the only
vertical divisors lie above base divisors. At general points of $Q$
and $T$, both divisors are smooth, $Y\to\bar Y$ is an isomorphism,
and the induced map $Q\to T$ is a submersion. A local equation of $T$
pulls back to $t^e$ times a unit. Take a holomorphic $e$th root of that
unit and absorb it into $t$; choose the remaining base coordinates
among coordinates on $Q$. This gives \eqref{eq:power-map}. Generic
smoothness and constructibility ensure that these neighborhoods occur
over a dense open of $T$, while avoiding any specified proper closed
subset of $Q$.
\end{proof}

A divisor dominating the base is called \emph{horizontal}; the others
are \emph{vertical}. The distinction in Lemma~\ref{lem:prepared-model}
will be used twice. Saturating pulled-back base differentials removes
the factor $t^{e-1}$ in \eqref{eq:power-map}; this is the ramification
correction in the slope proof. Later, the same coordinates estimate the
fiber integrals at base divisors. Exceptional divisors created on the
smooth resolution can map into higher codimension. They cause no loss
of divisor control because normalization is performed first on
$\bar Y$, and the final psh extensions are carried out on the smooth
base $S$.
